\subsubsection{Convex Hull (monotono)}

\paragraph{Idea intuitiva.}
Algoritmo de \textbf{Andrew's monotone chain}: ordenar puntos por $(x, y)$, luego construir la hull inferior y superior por separado. Para cada punto, aniadirlo al hull y, mientras el ultimo giro no sea ``counter-clockwise'' (\texttt{cross > 0}), popear. La demostracion: si los puntos se procesan en orden, el hull inferior y superior cubren la frontera completa sin auto-intersecciones.

\paragraph{Propiedades clave (invariantes).}
\begin{itemize}\setlength\itemsep{0pt}
	\item Ordena los puntos en $O(n \log n)$.
	\item Construye la hull en $O(n)$ despues.
	\item Output: poligono convexo en counter-clockwise, sin punto final duplicado.
	\item El \texttt{<= 0} incluye colineales; usar \texttt{< 0} para excluirlos.
	\item Cada punto aparece en la hull inferior o superior (no ambos).
\end{itemize}

\paragraph{Codigo clave.}
Construccion del hull inferior y superior:
\begin{verbatim}
sort(p.begin(), p.end());
vector<Point> lower, upper;
for (auto& p : points) {
    while (lower.size() >= 2 && cross(lower.back() - lower[lower.size()-2],
                                       p - lower.back()) <= 0)
        lower.pop_back();
    lower.push_back(p);
}
for (int i = points.size() - 1; i >= 0; --i) {
    auto& p = points[i];
    while (upper.size() >= 2 && cross(upper.back() - upper[upper.size()-2],
                                       p - upper.back()) <= 0)
        upper.pop_back();
    upper.push_back(p);
}
vector<Point> hull = lower;
hull.insert(hull.end(), upper.begin() + 1, upper.end() - 1);
\end{verbatim}

\paragraph{Complejidad.}
\begin{itemize}\setlength\itemsep{0pt}
	\item $O(n \log n)$.
	\item Constante muy pequena (solo cross products).
\end{itemize}

\paragraph{Donde tocar para modificar.}
\begin{itemize}\setlength\itemsep{0pt}
	\item \textbf{Colineales}: cambiar \texttt{<= 0} por \texttt{< 0} para excluirlos del hull.
	\item \textbf{Puntos repetidos}: \texttt{unique} despues de sort.
	\item \textbf{Hull 3D}: gift wrapping es $O(n^2)$; convex hull 3D es mas complejo.
	\item \textbf{Output con los puntos}: aniadir el hull completo (no solo upper/lower).
\end{itemize}

\paragraph{Trampas y errores frecuentes.}
\begin{itemize}\setlength\itemsep{0pt}
	\item El \texttt{pop\_back} y el \texttt{reverse} son esenciales para no duplicar el ultimo punto.
	\item Para excluir colineales, ajustar el \texttt{<= 0} vs \texttt{< 0}.
	\item Si los puntos estan duplicados, aniadir \texttt{unique} antes del sort.
\end{itemize}

\paragraph{Codigo.}
Ver \texttt{cpp.tex}, seccion \textit{Convex Hull (monotono)}.

\vspace{0.5em}
